% ===================================================================
%  TABELLA nr. 3 — Uffanza e giuventetgna.
%  Traduction 2026 d'apres texto/io, controlee sur texto/fr, texto/ca
%  et texto/oc. Consigne : voir texto/rm/10-tabella-01.tex.
%
%  LE CALENDRIER DE L'ALINEA 1 DE LA SECONDE SCENE PORTE LA TROUVAILLE
%  DE CETTE PAGE, ET ELLE TIENT EN DEUX MOTS.
%
%  Le romanche nomme les mois comme toutes les langues romanes —
%  « schaner », « favrer », « mars », « avrigl », « matg »,
%  « avust », « settember », « october », « november »,
%  « december » : ce sont les mois de Rome, et ils n'ont pas bouge.
%
%  MAIS DEUX D'ENTRE EUX ONT ETE REMPLACES, ET CE SONT PRECISEMENT
%  LES DEUX QUE CET ALINEA NOMME :
%    juin  : « zercladur », de SARCULARE, sarcler — LE MOIS OU L'ON
%            BINE ;
%    juillet : « fanadur », de FENUM, le foin — LE MOIS OU L'ON
%            FANE.
%  L'italien garde « giugno » et « luglio », le francais « juin » et
%  « juillet », tous de IUNIUS et de IULIUS. Le romanche a laisse
%  Auguste (« avust ») et a chasse Junon et Cesar.
%
%  POURQUOI CES DEUX-LA ? Parce que dans une vallee alpine ce sont
%  les deux seuls mois qui ont un TRAVAIL a eux, et un seul : on bine
%  en juin, on fane en juillet, et le reste de l'annee ne se laisse
%  pas resumer d'un geste. Les mois qu'on compte gardent le nom de
%  Rome ; les mois qu'on travaille prennent le nom du travail.
%
%  C'EST LA MEME LOI QUE LES TABLEAUX 15 LITUANIEN ET LUXEMBOURGEOIS
%  CHERCHAIENT SUR LA POMME DE TERRE, ET ELLE SE LIT ICI PLUS
%  NETTEMENT : ce n'est pas la date qui decide de la forme d'un nom,
%  c'est l'USAGE. Douze mois arrivent ensemble, par la meme route, a
%  la meme date ; deux seulement sont refaits, et ce sont les deux
%  qu'on a dans les mains.
%
%  ECART DE COMPOSITION DU BLOC t03-c2-01-1 : le fac-simile ido y
%  ecrit « junio (julio od agosto) » avec « julio » en romain, seul
%  des trois. gravuri/korekti.json remet le gras, comme pour les
%  vingt-et-une colonnes precedentes ; le fichier source, lui, ne
%  bouge pas.
% ===================================================================

%%K t03-apar-1 apar
\VUsaut{3.0mm}
\VUtitre{15.0pt}{60}{TABELLA nr. 3}
\VUsaut{1.6mm}
\VUtitre{11.4pt}{0}{Uffanza e giuventetgna.}
\VUsaut{3.4mm}

%%K t03-apar-2 apar
(Las tabellas 3 e 4 èn cumbinadas uschia ch'ellas preschentan en quatter
\VUgras{scenas da famiglia} las noziuns essenzialas davart l'\VUgras{aura},
l'\VUgras{età}, la \VUgras{famiglia}, la \VUgras{vestgadira} ed il \VUgras{stadi}.)

\VUsaut{4.0mm}
%%K t03-tit-1 sub
\VUcentre{12.0pt}{0}{\textit{Emprima scena.}}
\VUsaut{3.0mm}

%%K t03-tit-2 sub
\VUpk{10.2pt}{40}{Ceremonia da batten en in vitg.}
\VUsaut{2.4mm}

%%K t03-tit-3 sub
\VUpk{10.2pt}{40}{La primavaira.}
\VUsaut{3.0mm}

%%K t03-c1-01-1 p
1. --- Nus essan en la \VUgras{primavaira}, en ils mais \VUgras{mars}, \VUgras{avrigl} u
\VUgras{matg}, en ils dis da l'emna, \VUgras{glindesdi}, \VUgras{mardi} u
\VUgras{mesemna}. I è \VUgras{diesch uras} la damaun.

%%K t03-c1-02-1 p
2. --- L'\VUgras{aura} è bella, l'\VUgras{aria} pura. Il sulegl splendura. Intginas
\VUgras{nivulettas} \textsuperscript{(2)} van si en il blau \VUgras{tschiel}
\textsuperscript{(1)}.

%%K t03-c1-03-1 p
3. --- Ils \VUgras{arbers} han apaina \VUgras{fegls}. Ils sutgls \VUgras{papels}
\textsuperscript{(3)} ed il aut \VUgras{platan} \textsuperscript{(17)} han enfin ussa mo
buttuns.

%%K t03-c1-04-1 p
4. --- Las \VUgras{randulinas} \textsuperscript{(4)} returnan e sgolattan cun allegher
tschintschim enturn la \VUgras{frizza} \textsuperscript{(7)} dal \VUgras{clutger}
\textsuperscript{(5)}, che curunescha la \VUgras{baselgia} \textsuperscript{(6)} dal vitg.

%%K t03-tit-4 sub
\VUsaut{3.4mm}
\VUpk{10.2pt}{40}{La baselgia.}
\VUsaut{3.0mm}

%%K t03-c1-05-1 p
5. --- Fitg simpla è quella baselgia cun ses grond \VUgras{portal}
\textsuperscript{(13)} e sia gronda \VUgras{ura} \textsuperscript{(8)}. Il pitschen
\VUgras{punctader} \textsuperscript{(9)} è sin la cifra X, el mussa las \VUgras{uras}. Il
\VUgras{grond} \textsuperscript{(10)} è sin XII (mezdi); el mussa las \VUgras{minutas}.

%%K t03-c1-06-1 p
6. --- En il clutger sunan las \VUgras{campanas} \textsuperscript{(11)} allegramain. Sin
la culmegna dal tetg stat ina pitschna \VUgras{crusch} \textsuperscript{(12)} da crappa.

%%K t03-c1-07-1 p
7. --- La baselgia n'ha betg mo in grond \VUgras{portal} \textsuperscript{(13)}, ella ha
era ina purtella laterala. Tras questa purtella va ser \VUgras{plevon}
\textsuperscript{(15)}, vestgì cun sia \VUgras{sutana}, ord la \VUgras{sacrestia}
\textsuperscript{(14)} e vers la \VUgras{chasa parochiala} \textsuperscript{(19)}.

%%K t03-c1-08-1 p
8. --- Sper el, qua è la \VUgras{vendidra da latg} \textsuperscript{(24)} cun ses
\VUgras{asen} \textsuperscript{(23)}. Ella returna da la citad, nua ch'ella ha purtà bun
latg per ils uffants.

%%K t03-tit-5 sub
\VUsaut{4.0mm}
\VUpk{10.2pt}{40}{L'iert da fritga.}
\VUsaut{3.4mm}

%%K t03-c1-09-1 p
9. --- Segner Aliboron (l'asen) è cuntent da quella cursa da damaun. El aspira cun
delizia l'aria embalsamada dal parfum da las \VUgras{flurs} \textsuperscript{(20)}, che
cuvran ils \VUgras{arbers da fritga} \textsuperscript{(18)} da l'\VUgras{iert da fritga}
vischin. Tuts rosa ed alv èn quels \VUgras{persers} \textsuperscript{(18)} ed
\VUgras{abricosers} \textsuperscript{(19)}, ed els laschan sperar ina buna racolta.

%%K t03-c1-10-1 p
10. --- Entant van las pitschnas \VUgras{aviels} \textsuperscript{(22)} da la \VUgras{chasa
d'aviels} \textsuperscript{(21)} là a rimnar zumbrind lur dultsch \VUgras{mel}.

%%K t03-c1-11-1 p
11. --- Ma tge capita? Qua currin ils mats dal vitg e fan fugir las \VUgras{auchas}
\textsuperscript{(25)} spaventadas. Quai è la sortida ord la baselgia dal cortegi,
suenter il \VUgras{batten} dal figl dal notar.

%%K t03-c1-12-1 p
12. --- Il pitschen Jachen, en sia \VUgras{cuna} \textsuperscript{(26)}, sa dasda pervi
dal sunar allegher da las campanas, e sia sora Madlaina, che vul era vesair il
cortegi da batten, roga sia mamma da vegnir a la pectinar ed a la vestgir sin il
sulier da la chasa.

%%K t03-c1-13-1 p
13. --- La mamma consenta. Ella metta ses \VUgras{foulard} \textsuperscript{(28)}, ses
\VUgras{corsagi} \textsuperscript{(29)} e ses \VUgras{scussal} \textsuperscript{(30)}, ed
ella vegn a seser sin ina pitschna sutga avant la porta. Madlaina ha mo sia
\VUgras{sutgotta} \textsuperscript{(31)} e ses \VUgras{corset} \textsuperscript{(32)}.

%%K t03-c1-14-1 p
14. --- Tge fortunada ch'ella è! Ella emblida sias flurs preferidas, ses
\VUgras{garofels} \textsuperscript{(34)}, sin ils quals sa mettan las
\VUgras{farfaglias} \textsuperscript{(35)}, ses \VUgras{tulipans} \textsuperscript{(36)},
sias \VUgras{flurs da matg} \textsuperscript{(37)}, en \VUgras{vaschs}
\textsuperscript{(38)} sin il \VUgras{mir} \textsuperscript{(33)}, e sias
\VUgras{rosas} \textsuperscript{(39)}, che furman ina saiv uschè bella. Ella
contempla la ritga vestgadira da las dunnas dal cortegi.

%%K t03-c1-15-1 p
15. --- Ils mats dal vitg n'observan betg fitg la vestgadira. Els sa prescheschan
e sa dattan urts per survegnir \VUgras{bombons} \textsuperscript{(55)} u
\VUgras{munaida} \textsuperscript{(52)}. In d'els bragia \textsuperscript{(40)}.

%%K t03-c1-16-1 p
16. --- In auter è ì sin il pe dal \VUgras{chaun} \textsuperscript{(42)}, che fugia
guaind e fa fugir la \VUgras{giaglina} \textsuperscript{(43)} e ses \VUgras{pulschins}
\textsuperscript{(44)}.

%%K t03-tit-6 sub
\VUsaut{5.0mm}
\VUpk{10.2pt}{40}{Il cortegi da batten.}
\VUsaut{4.0mm}

%%K t03-c1-17-1 p
17. --- Avant il cortegi va la \VUgras{nudrizza} \textsuperscript{(45)}. Ella ha ina
bella \VUgras{cuffa} \textsuperscript{(46)} cun lungas bendas. Ella porta il
\VUgras{novnaschì} \textsuperscript{(47)}, tut vestgì cun \VUgras{pizzas}
\textsuperscript{(48)}.

%%K t03-c1-18-1 p
18. --- Suenter la nudrizza vegnan il \VUgras{padrin} \textsuperscript{(49)} e la
\VUgras{madrina} \textsuperscript{(53)}. Els spartan ils bombons e la munaida. Il
padrin ha \VUgras{guants} \textsuperscript{(51)} ed in \VUgras{chapè da tuba}
\textsuperscript{(50)}. La madrina tegna en maun in bel \VUgras{bugliet}
\textsuperscript{(54)}.

%%K t03-c1-19-1 p
19. --- Suenter els vesain nus l'\VUgras{aug} \textsuperscript{(56)} e l'\VUgras{onda}
\textsuperscript{(58)} da l'uffant. L'onda ha in elegant \VUgras{chapè}
\textsuperscript{(59)}, l'aug in nair \VUgras{redingot} \textsuperscript{(57)} ed in
gilet alv. Qua vegnan lura il \VUgras{cusrin} \textsuperscript{(60)} e la
\VUgras{cusrina} \textsuperscript{(61)}. Questa tegna en maun la \VUgras{sora}
\textsuperscript{(62)} dal novnaschì, la pitschna Berta.

%%K t03-c1-20-1 p
20. --- L'auter \VUgras{frar} \textsuperscript{(63)} da Berta va suenter. El dat il
maun ad ina \VUgras{parenta} \textsuperscript{(64)}. Ils \VUgras{amis}
\textsuperscript{(65)} da la famiglia vegnan lura.

%%K t03-tit-7 sub
\VUsaut{4.6mm}
\VUpk{10.2pt}{40}{Ils spectaturs.}
\VUsaut{3.4mm}

%%K t03-c1-21-1 p
21. --- Sper il cortegi, qua è in \VUgras{mendicant} \textsuperscript{(66)} sustegnì
sin sia \VUgras{crutscha} \textsuperscript{(67)}. El dumonda l'almosna.

%%K t03-c1-22-1 p
22. --- Da l'autra vart è in pitschen mat fitg \VUgras{polit}
\textsuperscript{(68)}. El salida e giavischa «\VUgras{bun di}» a las persunas
ch'el enconuscha.

%%K t03-c1-23-1 p
23. --- Ils vitgadins èn sortids da lur chasas per vesair il cortegi da batten. Nus
vesain in \VUgras{pur} \textsuperscript{(69)} cun ses \VUgras{bunet da cutun} e sper
el ina \VUgras{pura} \textsuperscript{(70)}. Lura ina \VUgras{veglia}
\textsuperscript{(71)} sustegnida sin ses \VUgras{bastun} \textsuperscript{(72)}. A la
fin il \VUgras{mulinar} \textsuperscript{(73)} cun ses grond \VUgras{chapè da feltr}
\textsuperscript{(74)} ed ina giuvna pura cun \VUgras{zoccels da lain}
\textsuperscript{(75)} als peis, che mussa a ses uffantin da chaminar.

%%K t03-c1-24-1 p
24. --- Quella scena è fitg divertenta.

\VUsaut{4.0mm}
%%K t03-tit-8 sub
\VUcentre{12.0pt}{0}{\textit{Segunda scena.}}
\VUsaut{3.0mm}

%%K t03-tit-9 sub
\VUpk{10.2pt}{40}{Festa publica.}
\VUsaut{2.4mm}

%%K t03-tit-10 sub
\VUpk{10.2pt}{40}{Gieus navals e regatas.}
\VUsaut{3.0mm}

%%K t03-c2-01-1 p
1. --- La \VUgras{stad} è arrivada. Il sulegl chaudont dal mais \VUgras{zercladur}
(fanadur u \VUgras{avust}) sfrunta ils giuvens ad ir a bagnar ed a far battè.

%%K t03-c2-02-1 p
2. --- Sin la riva dretga dal flum sa dasvestgian ils battellisits per prender
part als gieus navals ed a las regatas.

%%K t03-c2-03-1 p
3. --- In d'els trai giu ses \VUgras{chalzers} \textsuperscript{(1)}; el als
deslià. Ses dus cumpogns èn gia pronts. Uss han els mo lur \VUgras{camischola da
maglia} \textsuperscript{(2)} e lur \VUgras{chautschas da bagn} \textsuperscript{(3)}.
Els discurran spetgond lur amis. Els discurran da la fiera e da la festa che ha
lieu sin l'autra riva.

%%K t03-c2-04-1 p
4. --- Per terra, sper els, giaschan las \VUgras{vestgadiras} da lur cumpogns: in
\VUgras{pèr} da \VUgras{buttinas} \textsuperscript{(4)}, \VUgras{chalzers}
\textsuperscript{(5)}, \VUgras{chaltschettas} \textsuperscript{(6)}, chapels da
\VUgras{feltr} \textsuperscript{(7)} e da \VUgras{paglia} \textsuperscript{(8)} ed ina
\VUgras{cravatta} \textsuperscript{(9)}.

%%K t03-c2-05-1 p
5. --- In dals giuvenots ha plià premurusamain ses \VUgras{vestgì}
\textsuperscript{(10)}. El al metta sin ina pezza da toaletta, en la quala ses
vischin vegn prest a serrar en las duas \VUgras{vestgadiras}.

%%K t03-c2-06-1 p
6. --- Il sisavel mat resta enavos. El ha anc sia \VUgras{casquetta}
\textsuperscript{(11)} sin il chau. El n'ha tratg giu ni sia \VUgras{chamischa}
\textsuperscript{(12)}, ni ses \VUgras{collar} \textsuperscript{(13)}, ni sias
\VUgras{manschettas} \textsuperscript{(14)}. El tegna en maun ses \VUgras{gilet}
\textsuperscript{(17)} ed ha anc sias \VUgras{chautschas} \textsuperscript{(16)} e
sias \VUgras{bretellas} \textsuperscript{(15)}. Segir che el na vegn betg a esser
pront per la proxima cursa.

%%K t03-c2-07-1 p
7. --- Sper ils giuvenots maina in truoi, da la vart dal qual i ha blers
\VUgras{chardums} \textsuperscript{(18)}, a l'ur dal flum, nua ch'il bab Jachen
prepara tranter las \VUgras{channas} \textsuperscript{(19)} il \VUgras{fieu
d'artifizi} \textsuperscript{(20)} ch'ins vegn a trair la saira.

%%K t03-tit-11 sub
\VUsaut{4.4mm}
\VUpk{10.2pt}{40}{La cursa.}
\VUsaut{3.4mm}

%%K t03-c2-08-1 p
8. --- Sin il flum sa spassegian intginas dunnas en \VUgras{barca}
\textsuperscript{(21)}, sa protegend cun lur parasols. Ellas suondan la cursa, ch'è
fitg interessanta. En mintga \VUgras{iole} \textsuperscript{(22)} vogan quatter
\VUgras{vogaders} \textsuperscript{(24)} cun tut lur forzas, entant ch'il
\VUgras{timunier} \textsuperscript{(26)} tegna il \VUgras{timun}
\textsuperscript{(25)} e dirigia la fragila barca. Las \VUgras{vogas}
\textsuperscript{(23)} battan en cadenza l'aua e las lediras barcas van cun ina
sveltezza che fa vertigin.

%%K t03-c2-09-1 p
9. --- Intgins giuvenots luttan, sper il pilaster da la punt, per il premi dal
\VUgras{mast unct} \textsuperscript{(28)}. In d'els va prudentamain sin il mast
flexibel e sbrischaivel per cuntanscher la pitschna \VUgras{bandiera}
\textsuperscript{(29)} fixada a la fin. Ses \VUgras{concurrent}
\textsuperscript{(30)} malfortunà è crudà en l'aua. El nuda per vegnir a terra.

%%K t03-c2-10-1 p
10. --- Giuvenots pli vegls fan \VUgras{turnaments da barca} \textsuperscript{(32)}
sper il \VUgras{cai} \textsuperscript{(33)}. A tut quels gieus assista in numerus
\VUgras{public} \textsuperscript{(31)} sustegnì al \VUgras{parapet} da la
\VUgras{punt} \textsuperscript{(27)} ed amassà sin la \VUgras{riva}. Intgins
spectaturs sa volvan tuttina per suondar il gieu dal \VUgras{mast dal premi}
\textsuperscript{(45)} u per admirar il \VUgras{cortegi communal}, che vegn per la
\VUgras{spartida} dals premis.

%%K t03-c2-11-1 p
11. --- L'emprim, qua èn intgins \VUgras{mats} \textsuperscript{(35)} che van avant
ils \VUgras{musicants} \textsuperscript{(36)}; lura vegnan il \VUgras{sindic}
\textsuperscript{(37)}, ils assessurs ed il \VUgras{cussegl communal} da la
citadina. A la fin marschan ils \VUgras{pumpiers} \textsuperscript{(38)}.

%%K t03-tit-12 sub
\VUsaut{5.0mm}
\VUpk{10.2pt}{40}{La fiera.}
\VUsaut{3.6mm}

%%K t03-c2-12-1 p
12. --- Numerusa glieud è sin la \VUgras{plazza da fiera} \textsuperscript{(39)}. Ins
vegn a vesair las differentas \VUgras{barachas}: ils \VUgras{chavals da lain}
\textsuperscript{(40)}, il \VUgras{circus} \textsuperscript{(41)}, nua che la
represchentaziun è gia cumenzada; il \VUgras{bal public} \textsuperscript{(42)}, nua
ch'ils giuvenots e las giuvnas da la citad giaudan il plaschair dal
\VUgras{ballar}.

%%K t03-c2-13-1 p
13. --- En il fund vesain nus las \VUgras{arenas} \textsuperscript{(44)}, nua ch'il
valurus \VUgras{matador} spagnol lutta cunter il \VUgras{tor} \textsuperscript{(45)}
furius, lura la \VUgras{muntogna russa} \textsuperscript{(49)}, il \VUgras{tir}
\textsuperscript{(46)}, nua ch'ins sa exercitescha a duvrar las \VUgras{armas da
fieu} ed a trair sin il \VUgras{scopo}.

%%K t03-c2-14-1 p
14. --- Dasperas, qua è il \VUgras{teater da fiera} \textsuperscript{(47)}, nua ch'ins
gioga la \VUgras{cumedia} ed il \VUgras{drama}. Fitg dasperas è la baracha dal
\VUgras{gigant} \textsuperscript{(48)} e dal \VUgras{nanin}. La \VUgras{statura} dal
emprim è dus meters e quindesch centimeters; il segund è apaina uschè grond sco in
uffant.

%%K t03-c2-15-1 p
15. --- A la fin, qua è la gronda \VUgras{menagaria} \textsuperscript{(50)} cun sias
\VUgras{bes-chas selvaticas}, ses \VUgras{tiger} \textsuperscript{(51)} cun pussantas
\VUgras{unglas}, ses \VUgras{liun} \textsuperscript{(52)} cun densa \VUgras{criniera},
sias duas \VUgras{giraffas} \textsuperscript{(53)} e ses \VUgras{elefant}
\textsuperscript{(54)}, che duvra uschè abilmain sia \VUgras{trumba}.

%%K t03-c2-16-1 p
16. --- Il \VUgras{domadur} \textsuperscript{(55)}, che dumagna quellas bes-chas, è
a la porta da la baracha.

\VUsaut{6.0mm}
\VUfilet{22mm}
